% Part: second-order-logic
% Chapter: syntax-and-semantics
% Section: size-of-domain

\documentclass[../../../include/open-logic-section]{subfiles}

\begin{document}

\olfileid{sol}{syn}{siz}

\olsection{描述无穷与\usetoken{S}{enumerable}\usetoken{P}{domain}}

一个集合~$M$ 是（戴德金）无穷的，当且仅当存在一个单射函数~$f\colon M \to M$，它不是满射的，即满足 $\ran{f} \neq M$。
在一阶逻辑中，我们可以考虑一个一元函数符号~$f$，
并说在结构~$\Struct{M}$ 中指派给它的函数~$\Assign{f}{M}$ 是单射的且 $\ran{f} \neq \Domain{M}$：
\[
\lforall[x][\lforall[y][(\eq[f(x)][f(y)] \to \eq[x][y])]] \land
\lexists[y][\lforall[x][\eq/[y][f(x)]]].
\]
如果 $\Struct{M}$ 满足这个语句，$\Assign{f}{M}:
\Domain{M} \to \Domain{M}$ 是单射的，从而 $\Domain{M}$ 必定是无穷的。
如果 $\Domain{M}$ 是无穷的，因而这样的函数存在，我们可以令 $\Assign{f}{M}$ 为该函数，那么 $\Struct{M}$ 将满足该语句。
然而，这要求我们的语言包含我们用于此目的的非逻辑符号~$f$。
在二阶逻辑中，我们可以简单地断言这样的函数\emph{存在}。
这样就不需要~$f$了，从而得到纯二阶逻辑中的语句
\[
\fn{Inf} \ident \lexists[u][(\lforall[x][\lforall[y][(\eq[u(x)][u(y)]
      \lif \eq[x][y])]] \land \lexists[y][\lforall[x][\eq/[y][u(x)]]])].
\]
$\Sat{M}{\fn{Inf}}$ 当且仅当 $\Domain{M}$ 是无穷的。
于是我们可以定义 $\fn{Fin} \ident \lnot \fn{Inf}$；$\Sat{M}{\fn{Fin}}$ 当且仅当 $\Domain{M}$ 是有穷的。
没有任何一个纯一阶逻辑的单个语句能表达论域是无穷的，尽管一个无穷的语句集可以做到。
也不存在这样的纯一阶逻辑语句集：一个结构满足它当且仅当其论域是有穷的。

\begin{prop}
$\Sat{M}{\fn{Inf}}$ 当且仅当 $\Domain{M}$ 是无穷的。
\end{prop}

\begin{proof}
$\Sat{M}{\fn{Inf}}$ 当且仅当
  $\Sat{M}{\lforall[x][\lforall[y][(\eq[u(x)][u(y)] \lif \eq[x][y])]]
    \land \lexists[y][\lforall[x][\eq/[y][u(x)]]]}[s]$ 对
  某个变元指派~$s$ 成立。若它成立，则 $s(u)$ 是一个单射函数，并且存在某个 $y
  \in \Domain{M}$ 不在~$s(u)$ 的值域中。反之，如果存在一个单射 $f\colon \Domain{M} \to \Domain{M}$ 且其值域 $\ran{f}
  \neq \Domain{M}$，那么 $s(u) = f$ 就是这样一个变元
  指派。
\end{proof}

一个集合 $M$ 是可数的，如果存在其元素的一个枚举
\[
m_0, m_1, m_2, \dots
\]
（无重复但可能是有穷的）。这样的枚举存在，当且仅当存在一个元素 $z \in M$ 和一个函数 $f\colon M \to M$，
使得 $z$, $f(z)$, $f(f(z))$, …… 就是~$M$ 的所有元素。因为如果枚举存在，令 $z = m_0$ 并令 $f(m_k) = m_{k+1}$ （或者若 $m_k$ 是枚举的最后一个元素，则令 $f(m_k) = m_k$）
即为所需的元素和函数。另一方面，如果这样的 $z$ 和 $f$ 存在，那么 $z$, $f(z)$, $f(f(z))$, …… 就是~$M$ 的一个枚举，从而 $M$ 是可数的。我们可以用二阶逻辑来表达 $z$ 和~$f$ 的存在性，从而构造一个语句，它在一个结构中为真当且仅当该结构是可数的：
\[
\fn{Count} \ident
\lexists[z][\lexists[u][\lforall[X][((X(z) \land
      \lforall[x][(X(x) \lif X(u(x)))]) \lif \lforall[x][X(x)])]]]
\]

\begin{prop}
$\Sat{M}{\fn{Count}}$ 当且仅当 $\Domain{M}$ 是可数的。
\end{prop}

\begin{proof}
假设 $\Domain{M}$ 是可数的，并设 $m_0$, $m_1$, …… 是一个枚举。通过去除重复项，我们可以保证没有 $m_k$ 出现两次。定义 $f(m_k) = m_{k+1}$ 并令 $s(z) = m_0$ 和 $s(u) = f$。我们证明
\[
\Sat{M}{\lforall[X][((X(z) \land \lforall[x][(X(x) \lif X(u(x)))])
    \lif \lforall[x][X(x)])]}[s]
\]
取任意的 $M \subseteq \Domain{M}$。进一步假设
$\Sat{M}{(X(z) \land \lforall[x][(X(x) \lif
X(u(x)))])}[\Subst{s}{M}{X}]$。那么 $\Subst{s}{M}{X}(z) \in M$ 并且
只要 $x \in M$，就有 $(\Subst{s}{M}{X}(u))(x) \in M$。换句话说，由于 $\varAssign{\Subst{s}{M}{X}}{s}{X}$，$m_0 \in M$ 并且若
$x \in M$ 则 $f(x) \in M$，因此 $m_0 \in M$, $m_1 = f(m_0) \in M$,
$m_2 = f(f(m_0)) \in M$，等等。于是，$M = \Domain{M}$，所以
$\Sat{M}{\lforall[x][X(x)]}[\Subst{s}{M}{X}]$。由于 $M$ 是 $\Domain{M}$ 的任意子集，由此得证：
$\Sat{M}{\fn{Count}}$。

现在假设 $\Sat{M}{\fn{Count}}$，即
\[
\Sat{M}{\lforall[X][((X(z) \land \lforall[x][(X(x) \lif X(u(x)))])
    \lif \lforall[x][X(x)])]}[s]
\]
对某个变元指派~$s$ 成立。令 $m = s(z)$ 且 $f = s(u)$，并考虑 $M = \{m,
f(m), f(f(m)), \dots\}$。这样定义的 $M$ 显然是可数的。那么由假设，
\[
\Sat{M}{(X(z) \land \lforall[x][(X(x) \lif X(u(x)))])
    \lif \lforall[x][X(x)]}[\Subst{s}{M}{X}]
\]
此外，$\Sat{M}{X(z)}[\Subst{s}{M}{X}]$，因为 $M \ni m =
\Subst{s}{M}{X}(z)$；并且也有 $\Sat{M}{\lforall[x][(X(x) \lif
X(u(x)))]}[\Subst{s}{M}{X}]$，因为只要 $x \in M$ 就有 $f(x) \in
M$。所以，既然前件和条件句均被满足了，后件也必定被满足：$\Sat{M}{\lforall[x][X(x)]}[\Subst{s}{M}{X}]$。但这
意味着 $M = \Domain{M}$，而根据定义 $M$ 是可数的，因此 $\Domain{M}$ 是可数的。
\end{proof}

\begin{prob}
语句 $\fn{Inf} \land \fn{Count}$ 在且仅在可数无穷的论域中为真。调整 $\fn{Count}$ 的定义，使其成为另一个直接表达论域是可数无穷的不同语句，并证明它确实如此。
\end{prob}

\end{document}